\chapter{Discussion} \label{cha:discussion}

\section{Interpretation of Results}

The main research question asked how a Convolutional Neural Network compares to RS analysis in accurately and reliably detecting LSB steganography in 
PNG images. 
Taken as a whole, the results do not support a simple narrative that "deep learning is better". 
Both methods demonstrate distinct strengths, and the preferable method heavily depends on the specific embedding strategy and payload size being detected.

For Sequential and Random LSB embedding, the RS analysis performed better at every tested rate. 
At $r = 0.25$, RS showed 95.4\% accuracy on Sequential and 97.0\% on Random embedding, while the CNN scored 89.47\% on Sequential and 82.37\% on Random. 
The most significant difference between the two steganalysis methods is observed at low-payload embeddings (an advantage of roughly 15 percentage points 
in favor of RS), which gradually decreases at higher rates. 
These results perfectly align with the theoretical foundation of RS analysis, which states that Sequential and Random LSB replacement violate the 
natural balance between the number of pixels in $R_M$ and $R_{-M}$ groups, as established by \textcite{Fridrich2001}. 
This demonstrates that RS analysis is a highly optimized detector for the characteristic features of this specific image manipulation, 
and that generic CNNs offer no meaningful advantage over it in these cases.

The case for Edge Adaptive LSB, however, is very different. 
Both methods perform well at $r = 0.75$ (RS = 98.1\%, CNN = 94.73\%), but at $r = 0.25$, RS accuracy falls sharply to 72.5\% (AUC 0.761), while the CNN retains a 79.13\% accuracy. 
The adaptive technique's core mechanism—restricting modifications to textured, noisy regions—is the exact assumption violation that weakens RS analysis. 
Because the mathematical discrimination function already registers high values in those noisy regions, LSB flips no longer reliably push pixel 
groups from Regular to Singular, or vice versa. The CNN learns directly from the data distribution rather than relying on an explicit smoothness 
assumption, making it far less susceptible to this evasion tactic.

The crossover between the two methods on Edge Adaptive LSB embedding thus illustrates the theoretical argument made by \textcite{Qian2015} 
and \textcite{Apau2024}: the rigidity of handcrafted statistical features becomes a critical liability as embedding techniques begin to respect 
and mimic natural image statistics.

The most stark contrast can be observed in the tests on LSB Matching. 
Here, the RS analysis operates very closely to flipping a coin; with AUCs of 0.516, 0.504, and 0.514, it is virtually a random classifier. 
This is not a surprising outcome, as \textcite{Ker2005} observed that the $\pm 1$ adjustment in LSB Matching does not flip bits deterministically, 
but rather disturbs pixel values symmetrically. 
This process does not break the $R_M \approx R_{-M}$ symmetry assumption that RS analysis inherently relies upon. 
Conversely, the CNN accuracy ranged from 68.60\% at $r = 0.25$ to 91.37\% at $r = 0.75$. While 68.60\% accuracy is not flawless, 
it is significantly better than a random guess and demonstrates that the neural network can learn to identify artifacts that render RS analysis entirely blind. 
For the Matching LSB method, this comparison demonstrates which method is applicable at all, rather than merely which method performs better.

\section{Sub-Question: Sensitivity to Payload Size}

Our sub-question asked how detection performance varies as the hidden payload decreases. 
A degradation in performance can be observed across both methods, but the specific trajectory of that degradation is highly informative.

The degradation of RS analysis is heavily dependent on the embedding technique used to hide the payload. 
It performs exceptionally well for Sequential and Random embedding, as the drop in accuracy from $r = 0.75$ to $r = 0.25$ is only 
about four to five percentage points. 
Thus, RS analysis proves to be highly reliable even at the lowest tested embedding rates for these schemes.

In the case of Adaptive embedding, this same decrease in payload results in a much more severe drop in performance (from 98.1\% to 72.5\%). 
For Matching embedding, there is virtually no impact from varying the payload parameter, simply because the detector fails to identify the 
signal at any rate. 
This directly correlates with our findings discussed in Section 2.1.4, where the performance of statistical tests decreases when the induced 
asymmetry falls below the natural noise level of the image.

The CNN is much more consistent across degradation rates, exhibiting an average decrease of 7--14 percentage points from $r = 0.50$ to $r = 0.25$, 
and a smaller decrease from $r = 0.75$ to $r = 0.50$. Most importantly, there is no catastrophic failure in any embedding scheme. 
Even on Matching LSB, where RS fails completely at $r = 0.25$, the CNN maintains discriminative capacity. 
This supports the claim by \textcite{Apau2024} that learned representations generalize better to low-payload embeddings, even if, in this specific study, 
that advantage is only explicitly demonstrated against the two more advanced schemes.

Overall, the two approaches provide different answers to our sub-question depending on the embedding scheme. 
For Sequential and Random LSB, both approaches can be trusted to work down to $r = 0.25$, though RS is slightly more accurate. 
On Adaptive LSB, the CNN exhibits less catastrophic degradation, making it the preferred detector below $r = 0.50$. 
Finally, for Matching LSB, only the CNN shows any measurable degradation at all, as RS cannot distinguish the signals at any tested rate.

\section{Relation to Previous Work}

Three common assertions in existing steganalysis research are reaffirmed by these results. 
First, as established by \textcite{Fridrich2001} and \textcite{Fridrich2002}, RS analysis is a highly effective means of identifying classical 
LSB replacement, remaining robust across varying rates of embedded data. 
Second, as predicted by \textcite{Ker2005}, RS-type detectors fail against LSB Matching due to an inherent lack of structural compatibility with 
the embedding mechanics. 
Third, as suggested by \textcite{Qian2015}, \textcite{Wang2021}, and \textcite{Apau2024}, CNN-based steganalysis methods generalize much better 
across diverse embedding techniques, overcoming the limitations of rigid, handcrafted statistical tests.

This study confirms that the traditional steganalysis approach (RS) is highly successful within its intended operating range, 
maintaining superiority over neural approaches for basic spatial embeddings. 
Therefore, practitioners may confidently rely on RS analysis to detect plain LSB replacement in lossless PNGs, while reserving CNN methods for 
scenarios involving complex, adaptive, or matching-based embedding schemes.

\section{Limitations}

Numerous limitations apply to the conclusions drawn by this study, specifically regarding the reliance on a single cover source (BOSSBase 1.01) and a 
single image modality (grayscale images re-encoded as PNGs).
\textcite{Sedighi2016} found that the statistical characteristics of cover images significantly influence both embedding and detection performance. 
Therefore, the absolute accuracies obtained in this research should not be directly extended to other cover sources—such as consumer photographs or 
screen captures—without further empirical validation. 

Furthermore, the payload rates examined in this research are relatively high compared to contemporary steganalysis studies 
(which often evaluate payloads down to 0.1 bpp or lower). 
It would be highly relevant to observe how performance degradation occurs for these two methods when utilizing significantly smaller 
payloads (i.e., $< 0.05$ bpp), and to determine whether the relative advantage of CNNs becomes even more pronounced at these near-invisible 
payload sizes across varied cover sources.

\section{Future Work}

The limitations mentioned above naturally lead to three primary avenues for future research. 
The first avenue is to expand the range of payloads down to the sub-0.1 bpp range, allowing researchers to evaluate whether the degradation of 
CNNs remains graceful at extremely low embedding rates compared to statistical methods. 
The second avenue is to broaden the set of cover sources to include JPEGs decompressed to PNGs, smartphone photographs, and synthetically generated images. 
This would directly address the cover source dependency highlighted by \textcite{Sedighi2016} and provide a more stringent assessment of the CNN 
generalization hypothesis in real-world domains. 
The third avenue involves expanding the comparison to include highly modern, content-adaptive approaches such as S-UNIWARD and MiPOD. 
Aligning future comparisons with contemporary steganalysis benchmarking procedures will ultimately help determine whether traditional RS analysis 
retains any practical utility in modern contexts, or if the field must rely entirely on neural network detectors.

\section{Broader Impact}

The major implication of these findings is that the choice of methodology in applied steganalysis must be dictated by the specific threat model, 
rather than defaulting to the most "in vogue" deep learning method available. 
RS analysis remains a powerful, highly interpretable, and computationally efficient detector for plain LSB replacement. 
Its superior performance against low and high payloads in this analysis firmly endorses its continued use within forensic pipelines 
targeting basic steganographic content. 

However, when the target is an embedding created using an adaptive method or LSB Matching, a machine learning approach must be adopted. 
The CNN findings in this analysis demonstrate that deep learning can recover usable signals in contexts where RS analysis is completely blind, 
even if the model's accuracy falls slightly below the upper bounds reported in literature that utilizes massive computational resources and highly complex architectures. 
Ultimately, these two methods should not be viewed strictly as competitors, but as complementary tools. 
A highly viable real-world detection strategy would involve running both methods in parallel and ensembling the results to cover a broader 
spectrum of steganographic threats.

\section{Ethical Considerations}

These empirical findings allow for a nuanced ethical reflection. The results confirm that the simple, computationally inexpensive 
RS analysis remains highly sufficient against foundational techniques such as Sequential and Random LSB. 
This means that forensic practitioners and law enforcement do not necessarily need access to specialized, expensive hardware or 
massive training datasets to address a significant share of basic real-world cases. 

However, the results also highlight potential exploits by adversaries utilizing Matching LSB, which both RS analysis and our CNN struggled to 
detect reliably at low payloads ($r = 0.25$). 
We consider the disclosure of these vulnerabilities acceptable because the embedding strategies examined have been long-established in public 
literature \cite{Sharp2001, Ker2005, Luo2010}, and this study contributes no novel evasion methods. 
This research utilizes a standard academic dataset (BOSSBase 1.01) and strictly adheres to the principles of scientific integrity and defensive 
security research outlined in Chapter 3.